\section{Comunicazioni (802.11p, VLC, LTE/5G)}
\label{sec:sota_communications}

Le comunicazioni V2X per il platooning devono sostenere scambi periodici e a bassa latenza (beaconing e messaggi di controllo) in ambienti dinamici e spesso congestionati. IEEE~802.11p è stato a lungo il riferimento per le comunicazioni veicolari dedicate, con numerosi studi che ne analizzano prestazioni e limiti in scenari di platooning \cite{ieee80211p-2010,segata2015communication}. In parallelo, le tecnologie cellulari LTE/5G offrono nuove opportunità in termini di copertura e servizi, e sono oggi ampiamente studiate anche tramite simulatori di sistema integrabili in ambienti OMNeT++ \cite{nardini2020simu5g,3gpp-21914-1400}. Infine, la Visible Light Communication (VLC) è stata proposta come canale complementare per comunicazioni intraplatoon a corto raggio, con risultati promettenti in configurazioni ibride radio/VLC \cite{ishihara2015improving,schettler2019deeply}. In questa sezione si sintetizzano caratteristiche, punti di forza e \emph{failure modes} delle tecnologie considerate, preparando il terreno alla motivazione multi-canale.
